\chapter*{ВВЕДЕНИЕ}
\addcontentsline{toc}{chapter}{ВВЕДЕНИЕ}

В современном профессиональном спорте, и в частности в баскетболе (Национальная баскетбольная ассоциация --- NBA), сбор и анализ статистических данных играет ключевую роль. Если ранее оценка эффективности игроков базировалась преимущественно на базовых показателях (очки, подборы, передачи), то сегодня клубы, аналитики и скауты опираются на продвинутые метрики (Advanced Stats), такие как рейтинг эффективности игрока (PER), истинный процент попаданий (TS\%) или показатель полезности (BPM). 

Обработка таких данных требует надежных и производительных систем хранения, способных не только хранить огромные массивы исторических данных по каждому матчу, но и быстро агрегировать их по сезонам, командам и игрокам. В связи с этим, проектирование оптимизированной реляционной базы данных с механизмами кэширования является актуальной задачей.

\textbf{Целью курсовой работы} является разработка базы данных и программного обеспечения для анализа статистики игроков и команд NBA с реализацией расчетных метрик эффективности.

Для достижения поставленной цели необходимо решить следующие \textbf{задачи}:
\begin{enumerate}
	\item Провести анализ предметной области и существующих систем сбора спортивной статистики.
	\item Спроектировать структуру базы данных, нормализованную до третьей нормальной формы (3НФ), и определить ограничения целостности.
	\item Разработать алгоритмы расчета продвинутых метрик результативности на стороне СУБД.
	\item Разработать ролевую модель с разграничением прав доступа к данным.
	\item Спроектировать и реализовать клиент-серверное программное обеспечение для взаимодействия с базой данных.
	\item Интегрировать механизм кэширования результатов аналитических запросов.
	\item Провести исследование влияния стратегий индексации и кэширования на производительность системы при различных уровнях нагрузки.
\end{enumerate}

\textbf{Объектом исследования} выступают системы хранения и обработки спортивной статистики.

\textbf{Предметом исследования} являются методы проектирования реляционных баз данных, оптимизации аналитических запросов и кэширования данных.